\documentclass[12pt,letterpaper]{article}

\usepackage{fontspec}
\usepackage{polyglossia}
\setdefaultlanguage{spanish}
\setmainfont{Liberation Serif}
\setmonofont{Liberation Mono}

\usepackage{geometry}
\geometry{margin=2.5cm}

\usepackage{amsmath}
\usepackage{xcolor}
\usepackage{listings}
\usepackage{tikz}
\usetikzlibrary{positioning, arrows.meta, calc}

\tikzset{
    task/.style={
        rectangle, rounded corners, draw=codeblue!60, fill=blue!5,
        align=center, font=\small\ttfamily, inner sep=5pt,
        minimum height=0.8cm
    },
    wide/.style={
        task, text width=2.8cm, inner sep=3pt, font=\scriptsize\ttfamily
    },
    dep/.style={-{Latex[length=2mm]}, thick, draw=codeblue!70}
}
\usepackage{booktabs}
\usepackage{enumitem}
\usepackage{ragged2e}
\usepackage{hyperref}

\definecolor{codegray}{gray}{0.95}
\definecolor{codegreen}{rgb}{0.0,0.45,0.0}
\definecolor{codeblue}{rgb}{0.1,0.1,0.6}

\lstset{
    backgroundcolor=\color{codegray},
    basicstyle=\ttfamily\small,
    keywordstyle=\color{codeblue}\bfseries,
    commentstyle=\color{codegreen}\itshape,
    stringstyle=\color{red!70!black},
    numbers=left,
    numberstyle=\tiny\color{gray},
    numbersep=8pt,
    breaklines=true,
    showstringspaces=false,
    frame=single,
    rulecolor=\color{gray!50},
    tabsize=4,
    extendedchars=true
}

\title{Tarea 1. Modelado de soluciones paralelas}
\author{Islas G\'omez Luis Antonio \and Ram\'irez Candia Alfredo \and P\'erez Flores Julio C\'esar \and Navarro Soto Mario Alberto}
\date{\today}

\begin{document}

%% ------------------------------------------------------------------
%% PORTADA
%% ------------------------------------------------------------------
\begin{titlepage}
    \centering
    \vspace*{0.5cm}

    {\Large\bfseries BENEM\'ERITA UNIVERSIDAD AUT\'ONOMA DE PUEBLA\par}
    \vspace{0.4cm}
    {\large Facultad de Ciencias de la Computaci\'on\par}
    \vspace{1.2cm}

    {\large Programaci\'on Concurrente y Paralela\par}

    \vfill

    {\Huge\bfseries Tarea 1\par}
    \vspace{0.4cm}
    {\LARGE Modelado de soluciones paralelas\par}
    \vspace{0.3cm}
    {\large (Primer parcial)\par}

    \vfill

    {\large\bfseries Equipo:\par}
    \vspace{0.3cm}
    Islas G\'omez Luis Antonio --- 202372860\par
    Ram\'irez Candia Alfredo --- 202375947\par
    P\'erez Flores Julio C\'esar --- 202247445\par
    Navarro Soto Mario Alberto --- 202264602\par

    \vfill

    {\large Docente: Hilda Castillo Zacatelco\par}
    \vspace{0.6cm}
    {\large \today\par}
\end{titlepage}

\tableofcontents
\newpage

%% ------------------------------------------------------------------
\section{Introducci\'on y objetivo}
El objetivo de esta tarea es modelar una soluci\'on paralela para cada uno de
los ejercicios planteados, empleando tres herramientas de an\'alisis:
los \textbf{grafos de dependencia}, las sentencias \textbf{COBEGIN--COEND} y
las \textbf{condiciones de Bernstein}. El modelado no consiste en repartir el
trabajo ``a ojo'', sino en identificar qu\'e tareas pueden ejecutarse
simult\'aneamente y cu\'ales deben respetar un orden obligatorio.

%% ------------------------------------------------------------------
\section{Marco te\'orico}
\begin{itemize}[leftmargin=1.4cm]
    \item \textbf{COBEGIN--COEND:} notaci\'on de pseudoc\'odigo que expresa
    paralelismo (fork/join). Las sentencias encerradas entre \texttt{COBEGIN} y
    \texttt{COEND} son concurrentes: no existe un orden predefinido entre
    ellas, y el proceso padre espera a que \emph{todas} terminen antes de
    continuar despu\'es del \texttt{COEND}.

    \item \textbf{Condiciones de Bernstein:} dos bloques $S_i$ y $S_j$, con
    $L(S)$ el conjunto de variables que \emph{leen} y $E(S)$ el conjunto de
    variables que \emph{escriben}, pueden ejecutarse en paralelo si y solo si
    se cumplen las tres condiciones:
    \[
        L(S_i) \cap E(S_j) = \emptyset, \qquad
        E(S_i) \cap L(S_j) = \emptyset, \qquad
        E(S_i) \cap E(S_j) = \emptyset.
    \]
    Si se cumplen, los bloques son \emph{independientes} y van en el mismo
    \texttt{COBEGIN}; si alguna falla, existe una dependencia y deben colocarse
    en el orden que esta impone.

    \item \textbf{Grafo de dependencias:} los nodos son las tareas y una arista
    $A \rightarrow B$ indica que $B$ no puede comenzar hasta que $A$ termine.
    Los tipos cl\'asicos son RAW (dependencia verdadera: escribir y luego
    leer), WAR (antidependencia: leer y luego escribir) y WAW (dependencia de
    salida: escribir y luego escribir). Los nodos sin aristas entre s\'i
    forman el paralelismo disponible; el camino m\'as largo (camino cr\'itico)
    fija el tiempo m\'inimo y el ancho del grafo el grado de paralelismo.
\end{itemize}

%% ------------------------------------------------------------------
%% ------------------------------------------------------------------
%% Problema 1
%% ------------------------------------------------------------------
\section{Problema 1. Ordenamiento lexicogr\'afico}

\emph{Pendiente por desarrollar.}

% TODO: modelar el ordenamiento lexicogr\'afico ascendente de un arreglo de
% cadenas con grafos de dependencia, COBEGIN--COEND y condiciones de Bernstein.

%% ------------------------------------------------------------------
%% Problema 2
%% ------------------------------------------------------------------

% TODO: modelar la inversi\'on de un arreglo con grafos de dependencia,
% COBEGIN--COEND y condiciones de Bernstein.
\documentclass[12pt]{article}
\usepackage[utf8]{inputenc}
\usepackage[spanish]{babel}
\usepackage{geometry}
\geometry{letterpaper, margin=2.5cm}

% Paquetes de diseño y tablas
\usepackage{booktabs}
\usepackage{enumitem}

% Entorno para código y diagramas
\usepackage{listings}
\usepackage{xcolor}
\usepackage{tikz}
\usetikzlibrary{shapes.geometric, arrows.meta, positioning}

% Configuración de TikZ para los Grafos de Dependencia
\tikzset{
    task/.style={
        rectangle, 
        draw=black, 
        fill=gray!10, 
        thick, 
        rounded corners=2pt,
        align=center, 
        font=\small\sffamily,
        inner sep=6pt
    },
    dep/.style={
        ->, 
        >={Stealth[length=6pt]}, 
        thick, 
        shorten >=2pt, 
        shorten <=2pt
    }
}

% Configuración para bloques de seudocódigo
\lstset{
    basicstyle=\ttfamily\small,
    keywordstyle=\bfseries\color{blue!80!black},
    commentstyle=\itshape\color{green!50!black},
    frame=single,
    breaklines=true,
    showstringspaces=false
}

\begin{document}

% ==================================================================
% PORTADA INSTITUCIONAL (BUAP)
% ==================================================================
\begin{titlepage}
    \centering
    {\bfseries\Large BENEM\'ERITA UNIVERSIDAD AUT\'ONOMA DE PUEBLA \par}
    \vspace{0.3cm}
    {\large Facultad de Ciencias de la Computaci\'on \par}
    \vspace{1.5cm}
    {\large Programaci\'on Concurrente y Paralela \par}
    \vspace{2.5cm}
    
    {\bfseries\LARGE Tarea 1 \par}
    \vspace{0.4cm}
    {\Large Modelado de soluciones paralelas \par}
    \vspace{0.2cm}
    {\small (Primer parcial) \par}
    
    \vfill
    
    {\bfseries\large Equipo: \par}
    \vspace{0.3cm}
    \begin{tabular}{c}
        Islas G\'omez Luis Antonio --- 202372860 \\
        Ram\'irez Candia Alfredo --- 202375947 \\
        P\'erez Flores Julio C\'esar --- 202247445 \\
        Navarro Soto Mario Alberto --- 202264602 \\
    \end{tabular}
    
    \vfill
    
    {\large Docente: Hilda Castillo Zacatelco \par}
    \vspace{0.8cm}
    {\large 20 de septiembre de 2026 \par}
\end{titlepage}

\newpage
\tableofcontents
\newpage

% ==================================================================
% CONTENIDO DEL DOCUMENTO
% ==================================================================
\section{Introducci\'on}
El procesamiento paralelo permite descomponer problemas computacionales extensos en tareas independientes que se ejecutan simult\'aneamente en m\'ultiples unidades de procesamiento. En esta pr\'actica se formaliza el diseño de una soluci\'on paralela para la inversi\'on del orden de elementos en un arreglo unimensional.

\section{Objetivos}
\begin{itemize}
    \item Dise\~nar una soluci\'on paralela y eficiente para la inversi\'on de un arreglo de $N$ elementos.
    \item Evaluar las condiciones de Bernstein para validar la ausencia de condiciones de carrera en los accesos a memoria.
    \item Representar la precedencia y concurrencia de la soluci\'on mediante un Grafo de Dependencia.
    \item Estructurar el modelo concurrente haciendo uso del bloque expl\'icito \texttt{COBEGIN--COEND}.
\end{itemize}

\section{Marco Te\'orico}
El modelado de programas concurrentes exige mecanismos formales que garanticen la determinaci\'on de los resultados:
\begin{itemize}
    \item \textbf{Condiciones de Bernstein:} Establecen los requisitos de independencia de memoria entre dos tareas $S_i$ y $S_j$. Siendo $L(S)$ el conjunto de variables le\'idas y $E(S)$ el conjunto de variables escritas, se debe cumplir:
    \[
        L(S_i) \cap E(S_j) = \emptyset, \quad E(S_i) \cap L(S_j) = \emptyset, \quad E(S_i) \cap E(S_j) = \emptyset.
    \]
    \item \textbf{Grafos de Dependencia:} Es un grafo dirigido $G=(V, E)$ donde los v\'ertices $V$ representan las tareas operativas y las aristas dirigidas $E$ expresan dependencias de precedencia temporal.
    \item \textbf{COBEGIN--COEND:} Constructo sint\'actico para indicar la ejecuci\'on paralela de un conjunto de sentencias, garantizando una barrera de sincronizaci\'on impl\'icita al salir del bloque.
\end{itemize}

\section{Ejercicio 2: Inversi\'on de un Arreglo}

\subsection{Enunciado}
Dado un conjunto de n\'umeros en un arreglo $A$ de $N$ elementos, invertir el arreglo; es decir, el primer elemento deber\'a aparecer al final, el \'ultimo en la primera posici\'on, y as\'i sucesivamente.

\subsection{L\'ogica del Algoritmo Paralelo}
Para invertir el arreglo $A$, definimos $K = \lfloor N/2 \rfloor$ procesos concurrentes $S_0, S_1, \dots, S_{K-1}$. Cada proceso $S_i$ es responsable de intercambiar de manera independiente el elemento de la posici\'on $i$ con su contraparte en la posici\'on $N - 1 - i$:
\[
    S_i: \quad \text{intercambiar}(A[i], A[N - 1 - i]) \quad \text{para } 0 \le i < K.
\]
Si $N$ es impar, el elemento central $A[K]$ no requiere ning\'un movimiento.

\subsection{Condiciones de Bernstein}
Analizamos los conjuntos de Lectura ($L$) y Escritura ($E$) para cualquier par de procesos distintos $S_a$ y $S_b$ ($a \neq b$):
\begin{itemize}
    \item $L(S_a) = E(S_a) = \{ A[a], A[N - 1 - a] \}$
    \item $L(S_b) = E(S_b) = \{ A[b], A[N - 1 - b] \}$
\end{itemize}

Como $a \neq b$, las direcciones de memoria procesadas por $S_a$ y $S_b$ son totalmente disjuntas. Por lo tanto, se cumplen las tres condiciones de Bernstein:
\[
    L(S_a) \cap E(S_b) = \emptyset, \qquad 
    E(S_a) \cap L(S_b) = \emptyset, \qquad 
    E(S_a) \cap E(S_b) = \emptyset.
\]
Al no existir interferencia en memoria, los $K$ procesos pueden ejecutarse simult\'aneamente sin riesgos de condiciones de carrera.

\subsection{Grafo de Dependencia}
\begin{center}
\begin{tikzpicture}[node distance=1.2cm and 1.5cm]
    \node[task] (inicio) {Inicio};
    
    \node[task, below left=1.2cm and 2.0cm of inicio] (s0) {$S_0$: Swap($A[0], A[N-1]$)};
    \node[task, below left=1.2cm and 0.1cm of inicio] (s1) {$S_1$: Swap($A[1], A[N-2]$)};
    \node[below=1.4cm of inicio, font=\large] (dots) {$\cdots$};
    \node[task, below right=1.2cm and 2.0cm of inicio] (sk) {$S_{K-1}$: Swap($A[K-1], A[N-K]$)};
    
    \node[task, below=2.2cm of dots] (fin) {Fin};

    \draw[dep] (inicio) -- (s0);
    \draw[dep] (inicio) -- (s1);
    \draw[dep] (inicio) -- (sk);

    \draw[dep] (s0) -- (fin);
    \draw[dep] (s1) -- (fin);
    \draw[dep] (sk) -- (fin);
\end{tikzpicture}
\end{center}

\subsection{Implementaci\'on con COBEGIN--COEND}
\begin{lstlisting}
// Entrada: Arreglo A[N]
// K = N / 2

COBEGIN
    S_0;     // Intercambiar A[0] y A[N-1]
    S_1;     // Intercambiar A[1] y A[N-2]
    ...
    S_K_1;   // Intercambiar A[K-1] y A[N-K]
COEND

escribir A
\end{lstlisting}

\section{Conclusi\'on}
El intercambio simult\'aneo de pares sim\'etricos permite reducir la complejidad temporal de la inversi\'on de un arreglo de $\mathcal{O}(N)$ a $\mathcal{O}(1)$ en presencia de $N/2$ procesadores. El an\'alisis mediante las condiciones de Bernstein valida formalmente que esta descomposici\'on por pares es completamente determinista y libre de condiciones de carrera.

\end{document}

%% ------------------------------------------------------------------
%% Problema 3
%% ------------------------------------------------------------------
\section{Problema 3. Intersecci\'on de conjuntos}

\emph{Pendiente por desarrollar.}

% TODO: modelar el c\'alculo de la intersecci\'on de dos conjuntos de cadenas
% con grafos de dependencia, COBEGIN--COEND y condiciones de Bernstein.

%% ------------------------------------------------------------------
%% Problema 4
%% ------------------------------------------------------------------
\section{Problema 4. Ra\'ices reales de una ecuaci\'on cuadr\'atica}

\subsection*{Enunciado}
Determinar las ra\'ices reales de una ecuaci\'on cuadr\'atica
$ax^2 + bx + c = 0$ mediante la f\'ormula general.

\subsection*{An\'alisis del algoritmo}
El c\'alculo se descompone de la siguiente forma:
\[
    D = b^2 - 4ac, \qquad
    x_1 = \frac{-b + \sqrt{D}}{2a}, \qquad
    x_2 = \frac{-b - \sqrt{D}}{2a}.
\]
Si $D < 0$ no existen ra\'ices reales y se detiene el c\'alculo.

\subsection*{Condiciones de Bernstein}
Para cada tarea se identifican las variables que lee ($L$) y las que escribe
($E$):

\begin{center}
\begin{tabular}{lll}
\toprule
Tarea & Lee $L$ & Escribe $E$ \\
\midrule
T1: \texttt{b2 = b*b}          & $\{b\}$            & $\{b2\}$  \\
T2: \texttt{ac4 = 4*a*c}       & $\{a,c\}$          & $\{ac4\}$ \\
T3: \texttt{D = b2 - ac4}      & $\{b2,ac4\}$       & $\{D\}$   \\
T4: \texttt{r = sqrt(D)}       & $\{D\}$            & $\{r\}$   \\
T5: \texttt{x1 = (-b+r)/(2a)}  & $\{a,b,r\}$        & $\{x1\}$  \\
T6: \texttt{x2 = (-b-r)/(2a)}  & $\{a,b,r\}$        & $\{x2\}$  \\
\bottomrule
\end{tabular}
\end{center}

Verificaci\'on de independencia:
\begin{itemize}[leftmargin=1.4cm]
    \item \textbf{T1 $\parallel$ T2:} $L(\text{T1}) \cap E(\text{T2}) =
    \emptyset$, $E(\text{T1}) \cap L(\text{T2}) = \emptyset$ y
    $E(\text{T1}) \cap E(\text{T2}) = \emptyset$. Son independientes.
    \item \textbf{T5 $\parallel$ T6:} comparten \'unicamente \emph{lecturas}
    de $\{a,b,r\}$ y escriben variables distintas ($x1$ y $x2$), por lo que las
    tres condiciones se cumplen. Son independientes.
    \item \textbf{T1, T2 $\rightarrow$ T3 $\rightarrow$ T4:} existen
    dependencias RAW (T3 lee lo que T1 y T2 escriben; T4 lee lo que T3
    escribe). Deben respetar ese orden.
\end{itemize}

\subsection*{Grafo de dependencias}
El grafo tiene forma de ``diamante'': dos tareas iniciales en paralelo, un
punto de uni\'on y luego dos tareas finales en paralelo.

\begin{center}
\begin{tikzpicture}[node distance=1.2cm and 1.8cm]
    \node[task] (t1) {T1: b2 = b*b};
    \node[task, below=1.2cm of t1] (t2) {T2: ac4 = 4*a*c};
    \node[task, right=2cm of $(t1)!0.5!(t2)$] (t3) {T3: D = b2 - ac4};
    \node[task, right=1.8cm of t3] (t4) {T4: r = sqrt(D)};
    \node[task, right=1.8cm of t4, yshift=1.4cm] (t5) {T5: x1 = (-b+r)/(2a)};
    \node[task, right=1.8cm of t4, yshift=-1.4cm] (t6) {T6: x2 = (-b-r)/(2a)};
    \draw[dep] (t1) -- (t3);
    \draw[dep] (t2) -- (t3);
    \draw[dep] (t3) -- (t4);
    \draw[dep] (t4) -- (t5);
    \draw[dep] (t4) -- (t6);
\end{tikzpicture}
\end{center}

El camino cr\'itico es T1 $\rightarrow$ T3 $\rightarrow$ T4 $\rightarrow$ T5 y
el paralelismo m\'aximo es 2 (T1 con T2, y T5 con T6).

\subsection*{Soluci\'on con COBEGIN--COEND}
\begin{lstlisting}
leer a, b, c
COBEGIN
    b2  = b * b
    ac4 = 4 * a * c
COEND
D = b2 - ac4
si D < 0 entonces
    escribir "No hay raices reales"
sino
    r = sqrt(D)
    COBEGIN
        x1 = (-b + r) / (2*a)
        x2 = (-b - r) / (2*a)
    COEND
    escribir x1, x2
fin si
\end{lstlisting}

\noindent
Los dos primeros c\'alculos son independientes entre s\'i, por lo que van en el
primer \texttt{COBEGIN}. El c\'alculo de $D$ es el punto de uni\'on (join)
obligatorio. Una vez obtenida $r$, las dos ra\'ices se calculan en paralelo en
el segundo \texttt{COBEGIN}. La decisi\'on sobre $D < 0$ permanece secuencial
porque todas las tareas posteriores dependen de ella.

%% ------------------------------------------------------------------
%% Problema 5
%% ------------------------------------------------------------------
\section{Problema 5. M\'aximo de una matriz}

\subsection*{Enunciado}
Dado un arreglo bidimensional $A$ de tama\~no $N \times M$ ($N$ filas y $M$
columnas), se requiere dise\~nar un algoritmo concurrente que determine el valor
m\'aximo global de la matriz utilizando paralelismo de datos, modelado con
grafos de dependencia, condiciones de Bernstein y primitivas
\texttt{COBEGIN--COEND}.

\subsection*{An\'alisis del algoritmo}
La matriz se particiona horizontalmente asignando una fila completa a cada
proceso concurrente:
\begin{enumerate}[leftmargin=1.4cm]
    \item \textbf{Fase Concurrente ($S_0, S_1, \dots, S_{N-1}$):} Cada proceso
    $S_i$ calcula el valor m\'aximo de su fila correspondiente y lo almacena en
    la posici\'on $i$ de un vector auxiliar unidimensional $\text{max\_fila}$:
    \[
        \text{max\_fila}[i] = \max_{0 \le j < M} \{ A[i][j] \}.
    \]
    \item \textbf{Fase de Reducci\'on Secuencial ($S_{\text{final}}$):} Un proceso
    final recorre el arreglo auxiliar $\text{max\_fila}$ de longitud $N$ y
    extrae el m\'aximo global:
    \[
        \text{max\_global} = \max_{0 \le i < N} \{ \text{max\_fila}[i] \}.
    \]
\end{enumerate}

\subsection*{Condiciones de Bernstein}
Para cada proceso se identifican los conjuntos de variables que lee ($L$) y que
escribe ($E$):

\begin{center}
\begin{tabular}{lll}
\toprule
Proceso / Tarea & Lee $L$ & Escribe $E$ \\
\midrule
$S_i$ (fila $i$, con $0 \le i < N$) & $\{ A[i][j] \mid 0 \le j < M \}$ & $\{ \text{max\_fila}[i] \}$ \\
$S_{\text{final}}$ (reducci\'on)     & $\{ \text{max\_fila}[i] \mid 0 \le i < N \}$ & $\{ \text{max\_global} \}$ \\
\bottomrule
\end{tabular}
\end{center}

Verificaci\'on de independencia:
\begin{itemize}[leftmargin=1.4cm]
    \item \textbf{Entre filas independientes ($S_a \parallel S_b$, con $a \neq b$):}
    \[
        L(S_a) \cap E(S_b) = \emptyset, \qquad
        E(S_a) \cap L(S_b) = \emptyset, \qquad
        E(S_a) \cap E(S_b) = \emptyset.
    \]
    Cada proceso lee elementos de su propia fila y escribe en una posici\'on
    desacoplada del arreglo. Al cumplirse simult\'aneamente las tres
    condiciones, todos los procesos $\{S_0, S_1, \dots, S_{N-1}\}$ pueden
    ejecutarse en paralelo.
    \item \textbf{Dependencia con el proceso final ($S_i \rightarrow S_{\text{final}}$):}
    \[
        E(S_i) \cap L(S_{\text{final}}) = \{ \text{max\_fila}[i] \} \neq \emptyset.
    \]
    Existe una dependencia estricta de tipo RAW (\emph{Read-After-Write}). Por
    ende, $S_{\text{final}}$ no puede ejecutarse concurrentemente con las filas
    y requiere una barrera de sincronizaci\'on previa.
\end{itemize}

\subsection*{Grafo de dependencias}
Los procesos de fila se bifurcan en paralelo y convergen mediante una barrera
de sincronizaci\'on (\texttt{COEND}) antes de ejecutar la reducci\'on secuencial:

\begin{center}
\begin{tikzpicture}[node distance=1.4cm]
    \node[task] (inicio) {Inicio};
    \node[wide] (s0) at (-5.4,-2.0) {S0: max\_fila[0]};
    \node[wide] (s1) at (-1.8,-2.0) {S1: max\_fila[1]};
    \node[font=\large] (dots) at (1.8,-2.0) {$\cdots$};
    \node[wide] (sn) at (5.4,-2.0) {S\_N-1: max\_fila[N-1]};

    \node[wide] (sfinal) at (0,-3.8) {S\_final: max\_global};
    \node[task] (fin) at (0,-5.2) {Fin};

    \node[font=\scriptsize\itshape, text=red!70!black, align=center,
          anchor=west] at (1.9,-3.8)
          {Barrera de sincronizaci\'on\\(\texttt{COEND})};

    \draw[dep] (inicio) -- (s0);
    \draw[dep] (inicio) -- (s1);
    \draw[dep] (inicio) -- (sn);

    \draw[dep] (s0) -- (sfinal);
    \draw[dep] (s1) -- (sfinal);
    \draw[dep] (sn) -- (sfinal);

    \draw[dep] (sfinal) -- (fin);
\end{tikzpicture}
\end{center}

El camino cr\'itico es $\text{Inicio} \rightarrow S_i \rightarrow S_{\text{final}} \rightarrow \text{Fin}$ y el paralelismo m\'aximo es $N$.

\subsection*{Soluci\'on con COBEGIN--COEND}
\begin{lstlisting}
// Entrada: A[N][M] (Matriz de NxM)
// Auxiliar: max_fila[N]
// Salida: max_global

COBEGIN
    S_0;     // max_fila[0]   = max(A[0][0..M-1])
    S_1;     // max_fila[1]   = max(A[1][0..M-1])
    ...
    S_N_1;   // max_fila[N-1] = max(A[N-1][0..M-1])
COEND

S_final;     // max_global    = max(max_fila[0..N-1])
escribir max_global
\end{lstlisting}

\noindent
La b\'usqueda del m\'aximo en cada fila es totalmente independiente, por lo que
se ejecuta dentro del bloque \texttt{COBEGIN}. La variable \texttt{max\_global}
representa la consolidaci\'on de los resultados intermedios y s\'olo puede
calcularse una vez que la barrera \texttt{COEND} garantiza que todas las filas
concluyeron su trabajo.

%% ------------------------------------------------------------------
%% Problema 6
%% ------------------------------------------------------------------
\section{Problema 6. Alumnos aprobados en matem\'aticas}

\subsection*{Enunciado}
Mostrar el nombre de los alumnos que han aprobado el curso de matem\'aticas.
Para aprobar, \emph{cada} examen parcial debe tener calificaci\'on mayor o
igual que 6. Cada estudiante realiz\'o 3 ex\'amenes parciales que valen el 50\%
de su calificaci\'on y 4 actividades que valen el otro 50\%. De cada estudiante
se conoce el nombre, la edad y las calificaciones.

\subsection*{An\'alisis del algoritmo}
Para cada estudiante $i$ se calcula:
\[
    pp_i = \frac{p_{1i}+p_{2i}+p_{3i}}{3}, \qquad
    pa_i = \frac{a_{1i}+a_{2i}+a_{3i}+a_{4i}}{4},
\]
\[
    cf_i = 0.5\,pp_i + 0.5\,pa_i,
\]
\[
    \text{aprobado}_i =
    (p_{1i}\ge 6) \wedge (p_{2i}\ge 6) \wedge (p_{3i}\ge 6)
    \wedge (cf_i \ge 6).
\]

\subsection*{Condiciones de Bernstein}
El punto clave es que \textbf{los estudiantes son independientes entre s\'i}.
Para dos estudiantes $i \neq j$:
\[
    L(P_i) = \{p_{1i},p_{2i},p_{3i},a_{1i},\dots,a_{4i}\}, \quad
    E(P_i) = \{pp_i, pa_i, cf_i, \text{aprobado}_i\},
\]
y sus conjuntos son disjuntos de los del estudiante $j$. Por lo tanto
\[
    L(P_i)\cap E(P_j)=\emptyset,\qquad
    E(P_i)\cap L(P_j)=\emptyset,\qquad
    E(P_i)\cap E(P_j)=\emptyset,
\]
es decir, todos los estudiantes pueden procesarse en paralelo (paralelismo de
datos).

Dentro de cada estudiante, las tareas son:
\begin{center}
\begin{tabular}{lll}
\toprule
Tarea & Lee $L$ & Escribe $E$ \\
\midrule
T1: \texttt{pp = (p1+p2+p3)/3}              & $\{p_1,p_2,p_3\}$       & $\{pp\}$  \\
T2: \texttt{pa = (a1+a2+a3+a4)/4}           & $\{a_1,a_2,a_3,a_4\}$   & $\{pa\}$  \\
T3: \texttt{cf = 0.5*pp + 0.5*pa}           & $\{pp,pa\}$             & $\{cf\}$  \\
T4: \texttt{aprobado = (p1>=6)y...y(cf>=6)} & $\{p_1,p_2,p_3,cf\}$    & $\{\text{aprobado}\}$ \\
T5: imprimir \texttt{nombre} si aprobado    & $\{\text{aprobado},\text{nombre}\}$ & $\{\}$ \\
\bottomrule
\end{tabular}
\end{center}

\textbf{T1 $\parallel$ T2:} escriben variables distintas y leen datos distintos,
por lo que las tres condiciones de Bernstein se cumplen. \textbf{T3} es el
punto de uni\'on (depende de $pp$ y $pa$), \textbf{T4} depende de $cf$ y
\textbf{T5} de \texttt{aprobado}.

\subsection*{Grafo de dependencias}
Entre estudiantes no existe ninguna arista: el grafo queda formado por $N$
componentes independientes (paralelismo ideal). Dentro de cada estudiante s\'i
hay dependencias:

\begin{center}
\textbf{Entre estudiantes:} sin aristas ($N$ componentes independientes).

\vspace{0.4cm}
\begin{tikzpicture}[node distance=1.2cm]
    \node[task] (e1) {Estudiante 1};
    \node[task, right=1.2cm of e1] (e2) {Estudiante 2};
    \node[right=1.2cm of e2, font=\large] (dots) {$\cdots$};
    \node[task, right=1.2cm of dots] (en) {Estudiante N};
\end{tikzpicture}

\vspace{0.6cm}
\textbf{Dentro de cada estudiante:}

\vspace{0.4cm}
\begin{tikzpicture}[node distance=1.2cm and 1.6cm]
    \node[task] (t1) {T1: pp = (p1+p2+p3)/3};
    \node[task, below=1.2cm of t1] (t2) {T2: pa = (a1+a2+a3+a4)/4};
    \node[task, right=2cm of $(t1)!0.5!(t2)$] (t3) {T3: cf = 0.5*pp + 0.5*pa};
    \node[task, right=1.6cm of t3] (t4) {T4: aprobado = ... y (cf>=6)};
    \node[task, right=1.6cm of t4] (t5) {T5: imprimir nombre};
    \draw[dep] (t1) -- (t3);
    \draw[dep] (t2) -- (t3);
    \draw[dep] (t3) -- (t4);
    \draw[dep] (t4) -- (t5);
\end{tikzpicture}
\end{center}

\subsection*{Soluci\'on con COBEGIN--COEND}
\begin{lstlisting}
COBEGIN                         // un proceso por estudiante
    para i = 1 .. N
        COBEGIN
            pp[i] = (p1[i] + p2[i] + p3[i]) / 3
            pa[i] = (a1[i] + a2[i] + a3[i] + a4[i]) / 4
        COEND
        cf[i] = 0.5*pp[i] + 0.5*pa[i]
        aprobado[i] = (p1[i]>=6) y (p2[i]>=6) y (p3[i]>=6) y (cf[i]>=6)
COEND

para i = 1 .. N                 // impresion secuencial
    si aprobado[i] entonces escribir nombre[i]
\end{lstlisting}

\noindent
La impresi\'on es un \textbf{recurso compartido} (secci\'on cr\'itica): si se
realizara en paralelo habr\'ia que protegerla con un mutex para evitar salida
entrelazada; por eso se hace de forma secuencial. El c\'alculo, en cambio, es
totalmente paralelizable porque cada estudiante trabaja sobre sus propias
variables.

%% ------------------------------------------------------------------
%% Problema 7
%% ------------------------------------------------------------------
\section{Problema 7. Descuento en medicamentos}

% TODO: modelar la determinaci\'on del descuento en la compra de medicamentos
% con grafos de dependencia, COBEGIN--COEND y condiciones de Bernstein.
\documentclass[12pt]{article}
\usepackage[utf8]{inputenc}
\usepackage[spanish]{babel}
\usepackage{amsmath, amssymb, amsthm}
\usepackage{tikz}
\usetikzlibrary{arrows.meta, positioning, shapes.geometric}
\usepackage{geometry}
\geometry{letterpaper, margin=1in}

\begin{document}

\section{Problema 7. Redes de Interconexión y Topologías}

\subsection{Análisis Teórico: Hipercubos de Dimensión $n$}
Un hipercubo $n$-dimensional ($Q_n$) es un grafo cuyos vértices corresponden a las $2^n$ cadenas binarias de longitud $n$. Dos vértices están conectados por una arista si y solo si sus representaciones binarias difieren exactamente en un solo bit.

\subsubsection{Propiedades de la Topología}
\begin{itemize}
    \item \textbf{Número de Vértices (Nodos):} $V = 2^n$
    \item \textbf{Número de Aristas (Enlaces):} $E = n \cdot 2^{n-1}$
    \item \textbf{Grado de cada Nodo:} $D = n$
    \item \textbf{Diámetro de la Red:} $\text{Diámetro} = n$
\end{itemize}

\subsection{Representación Gráfica}
A continuación se ilustra la estructura de un hipercubo tridimensional ($Q_3$), conformado por $2^3 = 8$ nodos y $12$ enlaces binarios.

\begin{figure}[h!]
\centering
\begin{tikzpicture}[scale=2, auto, >=stealth,
    node distance=2cm,
    vertex/.style={circle, draw=blue!80, fill=blue!10, thick, inner sep=2pt, minimum size=8mm}]

    % Nodos del cubo frontal
    \node[vertex] (000) at (0,0) {000};
    \node[vertex] (100) at (2,0) {100};
    \node[vertex] (110) at (2,2) {110};
    \node[vertex] (010) at (0,2) {010};

    % Nodos del cubo trasero
    \node[vertex] (001) at (0.8,0.8) {001};
    \node[vertex] (101) at (2.8,0.8) {101};
    \node[vertex] (111) at (2.8,2.8) {111};
    \node[vertex] (011) at (0.8,2.8) {011};

    % Aristas del cubo frontal
    \draw[thick] (000) -- (100);
    \draw[thick] (100) -- (110);
    \draw[thick] (110) -- (010);
    \draw[thick] (010) -- (000);

    % Aristas del cubo trasero
    \draw[thick] (001) -- (101);
    \draw[thick] (101) -- (111);
    \draw[thick] (111) -- (011);
    \draw[thick] (011) -- (001);

    % Aristas de interconexión (frente a atrás)
    \draw[thick, dashed, red!70!black] (000) -- (001);
    \draw[thick, dashed, red!70!black] (100) -- (101);
    \draw[thick, dashed, red!70!black] (110) -- (111);
    \draw[thick, dashed, red!70!black] (010) -- (011);

\end{tikzpicture}
\caption{Topología de Hipercubo 3D ($Q_3$). Las líneas punteadas rojas unen nodos que difieren en el tercer bit.}
\end{figure}

\subsection{Conclusiones}
La topología en hipercubo ofrece un excelente compromiso entre el diámetro de la red y el grado de los nodos, permitiendo un enrutamiento eficiente de complejidad $\mathcal{O}(n)$ en algoritmos de paso de mensajes paralelos.

\end{document}


%% ------------------------------------------------------------------
\section{Conclusiones}
El modelado paralelo se apoya en el an\'alisis de dependencias y no en la
intuici\'on. Aplicando las condiciones de Bernstein se decide qu\'e bloques son
independientes y pueden ir dentro de un mismo \texttt{COBEGIN}, mientras que el
grafo de dependencias muestra el orden obligatorio y el paralelismo real
disponible. Los ejercicios con datos independientes por elemento (como el
problema 6) alcanzan un paralelismo casi perfecto, mientras que los que tienen
una cadena de c\'alculo com\'un (como el problema 4) quedan limitados por su
camino cr\'itico.

\end{document}
